\documentclass[11pt]{article}
\usepackage{fontspec,amsmath,amssymb,amsthm,hyperref}
\usepackage[margin=1in]{geometry}
\setlength{\emergencystretch}{3em}
\title{A reusable regular-interpolation interface}
\author{LangevinLean mathematical exposition}
\date{}
\begin{document}\maketitle
\noindent Source module: \texttt{UnadjustedLangevin/LSI/InterpolationData.lean}. Compile this proof document with XeLaTeX or LuaLaTeX. The accompanying HTML pairs the mathematical argument with the exact Lean statements and records the scope of the formalization.


The RegularInterpolationData structure bundles density paths identifying the actual ULA interpolation, regularity, weak-equation and error controls. Its methods derive path probability, finite KL, the one-step recurrence and the iterated entropy bound. Existence of this package is not asserted by its definition.


\section{Specify the actual laws, density regularity and energy data}


Package the relative densities \(r_{k,t}\), their time derivatives, and drift fields together with the exact law identity, positivity, spatial regularity, local entropy domination, finite score/error energies, admissible weak equation, drift-energy bound, and gradient-moment bound. These remain explicit mathematical inputs. No entropy derivative, entropy recurrence, or convergence estimate is a structure field.



\noindent Lean declarations: \texttt{RegularInterpolationData}.

\section{RegularInterpolationData.path probability}


The data's law field identifies the tilted density with the actual interpolation from \(\rho_k\). Probability preservation proves that interpolation is normalized. Rewrite by the law identity to obtain normalization of each density along the supplied path.



\noindent Lean declarations: \texttt{RegularInterpolationData.path\_probability}.

\section{RegularInterpolationData.kl finite}


Use local entropy regularity to prove finite KL for each relative density, then apply the time-zero law identification in \texttt{ula\_kl\_finite\_of\_interpolation}. The target probability instance is explicit; in the final theorem it is supplied by the LSI hypothesis.



\noindent Lean declarations: \texttt{RegularInterpolationData.kl\_finite}.

\section{RegularInterpolationData.one step}


Define \(H'_k(t)\) by the differentiated entropy integral. The regular weak-equation theorem proves its right-derivative identity and Fisher/drift production bound. Data fields supply probability, KL/Fisher finiteness, continuity, numerical error-energy and gradient bounds. Feed these proved ingredients to \texttt{ula\_kl\_one\_step\_of\_interpolation} to derive the recurrence rather than assuming it in the data.



\noindent Lean declarations: \texttt{RegularInterpolationData.one\_step}.

\section{RegularInterpolationData.kl bound}


Unpack every analytic field into \texttt{ula\_kl\_bound\_of\_regular\_interpolation}, retaining LSI, potential smoothness, and the step regime as separate arguments. Its already proved differentiation, one-step estimate, and geometric iteration give the full displayed KL bound for the actual ULA laws represented by the data.



\noindent Lean declarations: \texttt{RegularInterpolationData.kl\_bound}.

\end{document}
